\section{Classes and Objects}

\begin{frame}[fragile]{Classes and Objects}
    Every Python value is an \textbf{object}. A \textbf{class} is its blueprint.

    \begin{lstlisting}
class Vector:
    def __init__(self, x, y):
        self.x = x
        self.y = y

    def __str__(self):
        return f"Vector({self.x}, {self.y})"

    def __add__(self, other):
        return Vector(self.x + other.x,
                      self.y + other.y)

v1 = Vector(3, 4)
v2 = Vector(1, 2)
print(v1 + v2)   # Vector(4, 6)
    \end{lstlisting}
\end{frame}



\begin{frame}{Dunder Methods}
    \textbf{Magic (dunder) methods} let your class integrate with Python syntax.

    \vspace{0.5em}
    \begin{tabular}{lll}
    \toprule
    Method & Triggered by & PyTorch analogy \\
    \midrule
    \code{\_\_init\_\_}     & \code{MyClass(...)}   & layer definitions \\
    \code{\_\_repr\_\_}     & \code{repr(obj)}      & model summary string \\
    \code{\_\_call\_\_}     & \code{obj(x)}         & dispatches to \code{forward(x)} \\
    \code{\_\_len\_\_}      & \code{len(obj)}       & \code{Dataset.\_\_len\_\_} \\
    \code{\_\_getitem\_\_}  & \code{obj[i]}         & \code{Dataset.\_\_getitem\_\_} \\
    \code{\_\_add\_\_}      & \code{a + b}          & tensor operator overloading \\
    \bottomrule
    \end{tabular}

    \vspace{0.5em}
    \alert{Why now?} Every PyTorch model subclasses \code{torch.nn.Module} --- the same pattern.
\end{frame}

\begin{frame}[fragile]{Inheritance and \code{super()}}
    \begin{lstlisting}
class Module:
    def __call__(self, x):
        return self.forward(x)
    def forward(self, x):
        raise NotImplementedError

class ReLU(Module):
    def forward(self, x):
        return max(0.0, x)

class Sequential(Module):
    def __init__(self, *layers):
        self.layers = layers
    def forward(self, x):
        for layer in self.layers:
            x = layer(x)
        return x

model = Sequential(ReLU())
print(model(-1.0))  # 0.0
    \end{lstlisting}
\end{frame}

\begin{frame}[fragile]{Copy Objects}
    \alert{Common mistake:} assignment creates a \textbf{reference}, not a copy.

    \begin{columns}[t]
    \column{0.5\textwidth}
    \begin{block}{Shallow copy}
    \begin{lstlisting}
import copy
lst1 = [1, 2, [3, 4]]
lst2 = lst1.copy()

lst2[0] = 99     # lst1 unchanged
lst2[2][0] = 99  # lst1 also changes!
    \end{lstlisting}
    \end{block}

    \column{0.5\textwidth}
    \begin{block}{Deep copy}
    \begin{lstlisting}
import copy
lst1 = [1, 2, [3, 4]]
lst2 = copy.deepcopy(lst1)

lst2[2][0] = 99  # lst1 unchanged
    \end{lstlisting}
    \end{block}
    \end{columns}

    \vspace{0.5em}
    Use \code{copy.deepcopy()} whenever your object contains nested mutable structures.
\end{frame}

\begin{frame}{Copy Semantics in Memory}
\begin{center}
\scalebox{0.88}{%
\begin{tikzpicture}[
  font=\scriptsize,
  var/.style={draw=unistnavy, fill=unistemerald!30, rounded corners=2pt,
              font=\ttfamily\scriptsize, inner sep=3pt,
              minimum width=1.4cm, minimum height=0.45cm},
  obj/.style={draw=unistnavy, fill=gra, font=\ttfamily\scriptsize,
              inner sep=3pt, minimum width=2.1cm, minimum height=0.45cm},
  lbl/.style={font=\normalfont\bfseries\small, text=unistnavy},
  sub/.style={font=\ttfamily\tiny, text=gray},
  ptr/.style={-{Stealth[length=5pt]}, thick},
]

  %% Column dividers
  \draw[gray!35, dashed, thin] (-2.1, 2.9) -- (-2.1, -2.0);
  \draw[gray!35, dashed, thin] ( 2.1, 2.9) -- ( 2.1, -2.0);

  % ── 1. Reference assignment ───────────────────────────────────────────────
  \node[lbl] at (-4.2, 2.7) {Reference};
  \node[sub] at (-4.2, 2.3) {list2 = list1};

  \node[var] (r1) at (-4.9, 1.6) {list1};
  \node[var] (r2) at (-3.5, 1.6) {list2};

  \node[obj] (rA) at (-4.2, 0.5) {[\,1,\ 2,\ \textcolor{oge}{\textbullet}\,]};
  \node[obj] (rN) at (-4.2,-0.7) {[\,3,\ 4\,]};

  \draw[ptr, color=unistnavy] (r1.south) to[out=-90,in=150] (rA.north west);
  \draw[ptr, color=unistnavy] (r2.south) to[out=-90,in=30]  (rA.north east);
  \draw[ptr, color=oge]       (rA.south) -- (rN.north);

  \node[font=\normalfont\tiny, text=oge] at (-4.2,-1.4) {same object!};

  % ── 2. Shallow copy ───────────────────────────────────────────────────────
  \node[lbl] at (0, 2.7) {Shallow Copy};
  \node[sub] at (0, 2.3) {list2 = list1.copy()};

  \node[var] (s1) at (-0.75, 1.6) {list1};
  \node[var] (s2) at ( 0.75, 1.6) {list2};

  \node[obj] (sA) at (-0.75, 0.5) {[\,1,\ 2,\ \textcolor{oge}{\textbullet}\,]};
  \node[obj] (sB) at ( 0.75, 0.5) {[\,1,\ 2,\ \textcolor{oge}{\textbullet}\,]};
  \node[obj] (sN) at ( 0.0, -0.7) {[\,3,\ 4\,]};

  \draw[ptr, color=unistnavy] (s1.south) -- (sA.north);
  \draw[ptr, color=unistnavy] (s2.south) -- (sB.north);
  \draw[ptr, color=oge] (sA.south) to[out=-90,in=150] (sN.north west);
  \draw[ptr, color=oge] (sB.south) to[out=-90,in=30]  (sN.north east);

  \node[font=\normalfont\tiny, text=oge] at (0,-1.4) {shared nested!};

  % ── 3. Deep copy ──────────────────────────────────────────────────────────
  \node[lbl] at (4.2, 2.7) {Deep Copy};
  \node[sub] at (4.2, 2.3) {copy.deepcopy(list1)};

  \node[var] (d1) at (3.5, 1.6) {list1};
  \node[var] (d2) at (4.9, 1.6) {list2};

  \node[obj] (dA)  at (3.5, 0.5) {[\,1,\ 2,\ \textcolor{grn}{\textbullet}\,]};
  \node[obj] (dB)  at (4.9, 0.5) {[\,1,\ 2,\ \textcolor{grn}{\textbullet}\,]};
  \node[obj] (dN1) at (3.5,-0.7) {[\,3,\ 4\,]};
  \node[obj] (dN2) at (4.9,-0.7) {[\,3,\ 4\,]};

  \draw[ptr, color=unistnavy] (d1.south) -- (dA.north);
  \draw[ptr, color=unistnavy] (d2.south) -- (dB.north);
  \draw[ptr, color=grn] (dA.south) -- (dN1.north);
  \draw[ptr, color=grn] (dB.south) -- (dN2.north);

  \node[font=\normalfont\tiny, text=grn] at (4.2,-1.4) {fully independent};

\end{tikzpicture}}
\end{center}
\end{frame}

\begin{frame}[fragile]{Assert and Try-Except}
    \begin{columns}[t]
    \column{0.5\textwidth}
    \begin{block}{assert}
    \begin{lstlisting}
def add(a, b):
    return a + b

assert add(2, 3) == 5
assert add(-1, 1) == 0
print("PASSED!")
    \end{lstlisting}
    Raises \code{AssertionError} if \code{False}.
    \end{block}

    \column{0.5\textwidth}
    \begin{block}{try-except}
    \begin{lstlisting}
try:
    t = (1, 2, 3)
    t[0] = 10  # TypeError
except TypeError as e:
    print("Error:", e)
    \end{lstlisting}
    Catches exceptions gracefully.
    \end{block}
    \end{columns}
\end{frame}
